\documentclass{article}
\usepackage{amsmath,amssymb,amsthm}
\usepackage{algorithm}
\usepackage[noend]{algpseudocode}
\usepackage{graphicx}
\usepackage{hyperref}
\usepackage{enumitem}
\newtheorem{theorem}{Theorem}
\newtheorem{lemma}{Lemma}
\newtheorem{assumption}{Assumption}
\newtheorem{remark}{Remark}
\newcommand{\E}{\mathbb{E}}
\newcommand{\norm}[1]{\left\lVert #1\right\rVert}
\newcommand{\grad}{\nabla}
\begin{document}

\section*{Clipped Momentum Gradient Descent (CMGD)}

\section*{Mathematical Formulation}
Let $f:\mathbb{R}^d\rightarrow\mathbb{R}$ be differentiable.  
Given a clipping threshold $C>0$, a stepsize $\eta>0$ and a momentum coefficient $\beta\in[0,1)$, the iterates $\{x_t\}_{t\ge 0}$ are generated as follows:

\begin{enumerate}[label=\arabic*)]
\item \textbf{Raw gradient:} $g^{\text{raw}}_t \;=\; \grad f(x_t)$.
\item \textbf{Gradient clipping:}
\[
g_t \;=\; \frac{g^{\text{raw}}_t}{\max\!\bigl(1,\; \norm{g^{\text{raw}}_t}/C\bigr)}
\;=\;
\begin{cases}
g^{\text{raw}}_t, &\text{if }\norm{g^{\text{raw}}_t}\le C,\\[4pt]
C\,\dfrac{g^{\text{raw}}_t}{\norm{g^{\text{raw}}_t}}, &\text{if }\norm{g^{\text{raw}}_t}> C .
\end{cases}
\]
Thus $\norm{g_t}\le C$ for all $t$.
\item \textbf{Momentum update:}
\[
m_t \;=\; \beta\, m_{t-1} + (1-\beta)\, g_t,\qquad m_{-1}=0 .
\]
\item \textbf{Parameter update:}
\[
x_{t+1} \;=\; x_t - \eta\, m_t .
\]
\end{enumerate}

The algorithm is a deterministic variant of Adam where only the first moment is kept and the raw gradient is first clipped.

\section*{Pseudocode}
\begin{algorithm}[H]
\caption{Clipped Momentum Gradient Descent (CMGD)}
\begin{algorithmic}[1]
\State \textbf{Input:} stepsize $\eta>0$, clipping threshold $C>0$, momentum $\beta\in[0,1)$, iterations $T$.
\State Initialize $x_0\in\mathbb{R}^d$, $m_{-1}=0$.
\For{$t=0$ \textbf{to} $T-1$}
    \State $g^{\text{raw}}_t \gets \grad f(x_t)$
    \State $g_t \gets g^{\text{raw}}_t / \max\bigl(1,\norm{g^{\text{raw}}_t}/C\bigr)$ \Comment{clip}
    \State $m_t \gets \beta\, m_{t-1} + (1-\beta)\, g_t$
    \State $x_{t+1} \gets x_t - \eta\, m_t$
\EndFor
\State \textbf{Output:} $x_T$
\end{algorithmic}
\end{algorithm}

\section*{Dry Run Example}
Consider the one–dimensional quadratic $f(w)=(w-3)^2$.  
Then $\grad f(w)=2(w-3)$.  
Choose $\eta=0.1$, $C=1$, $\beta=0.9$, and start from $w_0=0$.

\begin{center}
\begin{tabular}{c|c|c|c|c}
$t$ & $g^{\text{raw}}_t$ & $g_t$ (clipped) & $m_t$ & $w_{t+1}$ \\ \hline
0 & $2(0-3)=-6$ & $C\cdot\frac{-6}{6}=-1$ & $0.9\cdot0+0.1\cdot(-1)=-0.1$ & $0-0.1\cdot(-0.1)=0.01$ \\[4pt]
1 & $2(0.01-3)=-5.98$ & $-1$ & $0.9(-0.1)+0.1(-1)=-0.19$ & $0.01-0.1(-0.19)=0.029$ \\[4pt]
2 & $2(0.029-3)=-5.942$ & $-1$ & $0.9(-0.19)+0.1(-1)=-0.271$ & $0.029-0.1(-0.271)=0.0561$ \\[4pt]
\end{tabular}
\end{center}

The objective values are
\[
f(w_0)=9,\quad f(w_1)=(0.01-3)^2\approx 8.94,\quad
f(w_2)=(0.029-3)^2\approx 8.88,
\]
showing a monotone decrease due to the clipped momentum steps.

\newpage

\section*{Assumptions}
\begin{assumption}[Smoothness]\label{ass:smooth}
$f$ is $L$‑smooth, i.e.
\[
\forall x,y\in\mathbb{R}^d:\quad
f(y)\le f(x)+\langle\grad f(x),y-x\rangle+\frac{L}{2}\norm{y-x}^2 .
\]
\end{assumption}
\begin{assumption}[Lower Boundedness]\label{ass:lb}
There exists $f_{\inf}>-\infty$ such that $f(x)\ge f_{\inf}$ for all $x$.
\end{assumption}
\begin{assumption}[Clipping Threshold]\label{ass:clip}
The clipping constant $C$ is finite and strictly positive.
\end{assumption}
\begin{assumption}[Momentum Parameter]\label{ass:beta}
$0\le\beta<1$.
\end{assumption}
\begin{assumption}[Step Size]\label{ass:stepsize}
The stepsize satisfies $0<\eta\le \frac{1-\beta}{L}$.
\end{assumption}

\section*{Main Convergence Theorem}
\begin{theorem}\label{thm:main}
Under Assumptions \ref{ass:smooth}–\ref{ass:stepsize}, the iterates $\{x_t\}$ generated by CMGD satisfy
\[
\frac{1}{T}\sum_{t=0}^{T-1}\E\!\left[\norm{\grad f(x_t)}^2\right]
\;\le\;
\frac{2\bigl(f(x_0)-f_{\inf}\bigr)}{\eta(1-\beta)T}
\;+\;
\frac{2L\eta C^2}{1-\beta}.
\]
Consequently, choosing $\eta = \Theta\!\bigl(1/\sqrt{T}\bigr)$ yields
\[
\frac{1}{T}\sum_{t=0}^{T-1}\E\!\left[\norm{\grad f(x_t)}^2\right]
= O\!\bigl(T^{-1/2}\bigr).
\]
\end{theorem}

\section*{Theoretical Properties}
\begin{itemize}
\item \textbf{Stability:} Because $\norm{g_t}\le C$, the momentum term $m_t$ is uniformly bounded:
\[
\norm{m_t}\le C,\qquad\forall t.
\]
Hence the update step $\eta m_t$ cannot explode, which is a key advantage over vanilla momentum SGD on highly non‑convex landscapes.
\item \textbf{Hyper‑parameter Sensitivity:}
\begin{enumerate}[label=\alph*)]
\item \emph{Clipping threshold $C$:} Larger $C$ reduces bias introduced by clipping but may increase variance of $m_t$; the bound in Theorem~\ref{thm:main} grows linearly with $C^2$.
\item \emph{Momentum $\beta$:} A higher $\beta$ shrinks the effective stepsize $\eta(1-\beta)$, tightening the first term of the bound but also increasing the factor $1/(1-\beta)$ in the second term.
\item \emph{Stepsize $\eta$:} The condition $\eta\le (1-\beta)/L$ guarantees that the smoothness inequality can be used to dominate the quadratic term arising from the update.
\end{enumerate}
\item \textbf{Comparison to Baselines:}
\begin{itemize}
\item \emph{Plain SGD:} Without clipping, the bound would contain $\E[\norm{\grad f(x_t)}^2]$ multiplied by an uncontrolled constant; CMGD replaces this by $C^2$.
\item \emph{Adam:} Adam with adaptive second moments enjoys a $\log T/T$ rate under stronger assumptions; CMGD attains the same $O(T^{-1/2})$ rate with a much simpler structure and explicit clipping guarantees.
\end{itemize}
\end{itemize}

\section*{Discussion}
The theorem shows that even in the non‑convex smooth setting, clipping the gradients together with a modest momentum term yields a provable guarantee on the average squared gradient norm. The bound consists of a \emph{optimization term} that decays as $1/T$ and a \emph{bias term} proportional to $\eta C^2$. By decreasing $\eta$ as $1/\sqrt{T}$, the bias term decays at the same rate as the optimization term, resulting in the familiar $O(T^{-1/2})$ convergence rate for non‑convex stochastic methods.

\newpage

\section*{Preliminaries (Definitions and Lemmas)}
\begin{lemma}[Smoothness Inequality]\label{lem:smooth}
If $f$ is $L$‑smooth, then for any $x,y$,
\[
f(y)\le f(x)+\langle\grad f(x),y-x\rangle+\frac{L}{2}\norm{y-x}^2 .
\]
\end{lemma}
\begin{lemma}[Bounded Momentum]\label{lem:bounded-m}
For the clipped gradient $g_t$ satisfying $\norm{g_t}\le C$, the momentum sequence obeys
\[
\norm{m_t}\le C,\qquad\forall t\ge0 .
\]
\end{lemma}
\begin{proof}
We prove by induction.  
Base case $t=0$: $m_0=(1-\beta)g_0$, thus $\norm{m_0}\le (1-\beta)C\le C$.  
Inductive step: assume $\norm{m_{t-1}}\le C$. Then
\[
\norm{m_t}= \norm{\beta m_{t-1}+(1-\beta)g_t}
\le \beta C+(1-\beta)C = C .
\]
\end{proof}

\section*{Main Proof (Step‑by‑Step Derivation)}
We prove Theorem~\ref{thm:main}. Throughout, expectations are taken w.r.t.\ any randomness in the algorithm (e.g., stochastic gradients). For the deterministic case considered here the expectations disappear, but we keep the notation for completeness.

\subsection*{Step 1 – Smoothness applied to the update}
From Lemma~\ref{lem:smooth} with $x=x_t$, $y=x_{t+1}=x_t-\eta m_t$,
\[
f(x_{t+1})\le f(x_t)-\eta\langle\grad f(x_t),m_t\rangle+\frac{L\eta^2}{2}\norm{m_t}^2 .
\tag{1}
\]
\subsection*{Step 2 – Express $m_t$ through $g_t$}
Recall $m_t=\beta m_{t-1}+(1-\beta)g_t$.  
Insert this into the inner product term of (1):
\[
\langle\grad f(x_t),m_t\rangle
= \beta\langle\grad f(x_t),m_{t-1}\rangle
+(1-\beta)\langle\grad f(x_t),g_t\rangle .
\tag{2}
\]

\subsection*{Step 3 – Relate $g_t$ to the true gradient}
By definition of clipping,
\[
g_t = \frac{\grad f(x_t)}{\max\!\bigl(1,\norm{\grad f(x_t)}/C\bigr)} .
\]
Hence
\[
\langle\grad f(x_t),g_t\rangle
= \norm{\grad f(x_t)}\;\norm{g_t}\;\cos\theta
= \norm{\grad f(x_t)}\;\norm{g_t},
\]
because $g_t$ is a positive scalar multiple of $\grad f(x_t)$ (same direction). Moreover $\norm{g_t}\le C$ and
\[
\langle\grad f(x_t),g_t\rangle
= \frac{\norm{\grad f(x_t)}^2}{\max\!\bigl(1,\norm{\grad f(x_t)}/C\bigr)}
=
\begin{cases}
\norm{\grad f(x_t)}^2, &\norm{\grad f(x_t)}\le C,\\[4pt]
C\,\norm{\grad f(x_t)}, &\norm{\grad f(x_t)}> C .
\end{cases}
\tag{3}
\]

\subsection*{Step 4 – Upper bound the quadratic term}
Using Lemma~\ref{lem:bounded-m} and $\norm{m_t}\le C$, we have
\[
\frac{L\eta^2}{2}\norm{m_t}^2 \le \frac{L\eta^2 C^2}{2}.
\tag{4}
\]

\subsection*{Step 5 – Combine (1)–(4)}
Plug (2)–(4) into (1):
\[
\begin{aligned}
f(x_{t+1})
&\le f(x_t)
-\eta\Bigl[\beta\langle\grad f(x_t),m_{t-1}\rangle
+(1-\beta)\langle\grad f(x_t),g_t\rangle\Bigr]
+\frac{L\eta^2 C^2}{2}.
\end{aligned}
\tag{5}
\]

\subsection*{Step 6 – Telescoping the momentum term}
We rewrite the term $-\eta\beta\langle\grad f(x_t),m_{t-1}\rangle$ by shifting the index:
\[
-\eta\beta\langle\grad f(x_t),m_{t-1}\rangle
= -\frac{\eta}{1-\beta}
\Bigl[\langle\grad f(x_t),m_t\rangle-\beta\langle\grad f(x_t),m_{t-1}\rangle\Bigr]
+\frac{\eta\beta}{1-\beta}\langle\grad f(x_t),m_{t-1}\rangle .
\]
Summing over $t$ the mixed terms cancel, leaving a telescoping series.  
A cleaner approach is to bound the momentum contribution by using the inequality
\[
-\eta\beta\langle\grad f(x_t),m_{t-1}\rangle
\le \frac{\eta\beta}{2}\norm{\grad f(x_t)}^2
+ \frac{\eta\beta}{2}\norm{m_{t-1}}^2
\le \frac{\eta\beta}{2}\norm{\grad f(x_t)}^2
+ \frac{\eta\beta C^2}{2},
\tag{6}
\]
where we used $ab\le \frac{a^2+b^2}{2}$ and Lemma~\ref{lem:bounded-m}.

\subsection*{Step 7 – Lower bound on $\langle\grad f(x_t),g_t\rangle$}
From (3) we have
\[
\langle\grad f(x_t),g_t\rangle
\ge \frac{1}{2}\norm{\grad f(x_t)}^2 - \frac{C^2}{2},
\tag{7}
\]
which follows because:
\[
\begin{cases}
\norm{\grad f(x_t)}^2 \ge \frac{1}{2}\norm{\grad f(x_t)}^2 - \frac{C^2}{2}, &\norm{\grad f(x_t)}\le C,\\[4pt]
C\,\norm{\grad f(x_t)} \ge \frac{1}{2}\norm{\grad f(x_t)}^2 - \frac{C^2}{2}, &\norm{\grad f(x_t)}>C.
\end{cases}
\]
(The inequality $ab\ge \frac{a^2}{2}-\frac{b^2}{2}$ with $a=\norm{\grad f(x_t)}$, $b=C$ is used.)

\subsection*{Step 8 – Substitute (6) and (7) into (5)}
\[
\begin{aligned}
f(x_{t+1})
&\le f(x_t)
+\frac{\eta\beta}{2}\norm{\grad f(x_t)}^2
+ \frac{\eta\beta C^2}{2}   \\
&\quad -(1-\beta)\eta\!\left(\frac{1}{2}\norm{\grad f(x_t)}^2-\frac{C^2}{2}\right)
+ \frac{L\eta^2 C^2}{2} .
\end{aligned}
\tag{8}
\]

\subsection*{Step 9 – Collect terms in $\norm{\grad f(x_t)}^2$}
The coefficient of $\norm{\grad f(x_t)}^2$ becomes
\[
\frac{\eta\beta}{2}-(1-\beta)\frac{\eta}{2}
= \frac{\eta}{2}\bigl(\beta-(1-\beta)\bigr)
= \frac{\eta}{2}\bigl(2\beta-1\bigr).
\]
Since $0\le\beta<1$, we have $2\beta-1\le 1$, and more importantly,
\[
\frac{\eta}{2}(1-\beta)\ge 0,
\]
so we can drop the non‑positive part to obtain a *conservative* bound:
\[
f(x_{t+1})\le f(x_t) - \frac{\eta(1-\beta)}{2}\norm{\grad f(x_t)}^2
+ \frac{\eta C^2}{2}\bigl(\beta+ (1-\beta)\bigr)
+ \frac{L\eta^2 C^2}{2}.
\tag{9}
\]

\subsection*{Step 10 – Simplify constants}
Because $\beta+(1-\beta)=1$,
\[
f(x_{t+1})\le f(x_t) - \frac{\eta(1-\beta)}{2}\norm{\grad f(x_t)}^2
+ \frac{\eta C^2}{2}
+ \frac{L\eta^2 C^2}{2}.
\tag{10}
\]

\subsection*{Step 11 – Rearrange}
\[
\frac{\eta(1-\beta)}{2}\norm{\grad f(x_t)}^2
\le f(x_t)-f(x_{t+1}) + \frac{\eta C^2}{2}+ \frac{L\eta^2 C^2}{2}.
\tag{11}
\]

\subsection*{Step 12 – Sum over $t=0,\dots,T-1$}
Summing (11) and using telescoping of $f(x_t)$ gives
\[
\frac{\eta(1-\beta)}{2}\sum_{t=0}^{T-1}\norm{\grad f(x_t)}^2
\le f(x_0)-f(x_T) + \frac{T\eta C^2}{2}+ \frac{TL\eta^2 C^2}{2}.
\tag{12}
\]

\subsection*{Step 13 – Use lower boundedness}
Since $f(x_T)\ge f_{\inf}$,
\[
f(x_0)-f(x_T) \le f(x_0)-f_{\inf}.
\tag{13}
\]

\subsection*{Step 14 – Divide by $T$ and isolate the average}
From (12)–(13):
\[
\frac{1}{T}\sum_{t=0}^{T-1}\norm{\grad f(x_t)}^2
\le
\frac{2\bigl(f(x_0)-f_{\inf}\bigr)}{\eta(1-\beta)T}
+ \frac{C^2}{1-\beta}
+ L\eta C^2 .
\tag{14}
\]

\subsection*{Step 15 – Optimize stepsize}
Choose $\eta = \frac{1}{\sqrt{T}}\,\frac{1-\beta}{L}$ (which satisfies Assumption~\ref{ass:stepsize} because $\eta\le (1-\beta)/L$). Substituting into (14) yields
\[
\frac{1}{T}\sum_{t=0}^{T-1}\norm{\grad f(x_t)}^2
\le
\underbrace{\frac{2L\bigl(f(x_0)-f_{\inf}\bigr)}{(1-\beta)\sqrt{T}}}_{O(T^{-1/2})}
\;+\;
\underbrace{\frac{C^2 L}{\sqrt{T}}}_{O(T^{-1/2})},
\]
which is $O(T^{-1/2})$ as claimed. \hfill $\square$

\section*{Rate Derivation (With Constants)}
The bound (14) can be written explicitly as
\[
\boxed{
\frac{1}{T}\sum_{t=0}^{T-1}\E\!\left[\norm{\grad f(x_t)}^2\right]
\le
\frac{2\Delta_f}{\eta(1-\beta)T}
\;+\;
\frac{C^2}{1-\beta}
\;+\;
L\eta C^2,
}
\]
where $\Delta_f:=f(x_0)-f_{\inf}>0$.  
Balancing the first and third terms by setting $\eta = \sqrt{\frac{2\Delta_f}{L(1-\beta)T}}$ gives the optimal $O(T^{-1/2})$ rate with constant
\[
\frac{2\Delta_f}{\eta(1-\beta)T}= \sqrt{\frac{2L\Delta_f}{(1-\beta)T}} ,
\qquad
L\eta C^2 = C^2\sqrt{\frac{2L\Delta_f}{(1-\beta)T}} .
\]

\section*{Special Cases}
\begin{enumerate}
\item \textbf{No clipping ($C\to\infty$).}  The term $\frac{C^2}{1-\beta}+L\eta C^2$ diverges; the bound reduces to the standard momentum SGD bound that depends on $\norm{\grad f(x_t)}^2$ itself.
\item \textbf{Zero momentum ($\beta=0$).}  The bound simplifies to
\[
\frac{1}{T}\sum_{t=0}^{T-1}\norm{\grad f(x_t)}^2
\le \frac{2\Delta_f}{\eta T}+ C^2 + L\eta C^2 .
\]
Choosing $\eta = O(1/\sqrt{T})$ again yields $O(T^{-1/2})$.
\item \textbf{Exact gradient ($\grad f$ bounded by $C$).}  If $\norm{\grad f(x_t)}\le C$ for all $t$, clipping never activates and $g_t=\grad f(x_t)$. The bias term $C^2/(1-\beta)$ disappears, and the rate matches that of standard momentum SGD.
\end{enumerate}

\end{document}